% Pursuing Stacks (cotes 134-2 to 134-8): archive cote 134-6, site page 78 (PDF page 79), typescript page 396.
% typescript. Draft compared: ps5.tex:5403-5493.
% AI-transcribed by Claude Opus 5.5 (claude-opus-5-5, adaptive thinking, effort high), Message Batches API, batch ps3 (msgbatch_01XdFmuVVBAcgc7EFxdtCqsj), 2026-11-12; image review recorded in records/CHECK_DECISIONS_134_3_4.tsv.
\scanpage{134-6/78}
\mspage{396}

the category of \hand{$\mathcal{M}$}-val\unsure{\over[typed]{a}{u}}\over[typed]{e}{e}d \uline{sheaves} on the topos $B^{\wedge}$.

This equivalence, defined up to unique isomorphism, is deduced
from the diagram of canonical equivalences
% Hand-drawn arrows are marked blue; typed objects remain plain.
\diagram{\begin{tikzcd}[column sep=small]
\underline{\mathrm{Hom}}_{\hand{!}}(A^{\wedge},\hand{\mathcal{M}})\hand{\overset{\mathrm{def}}{=}(A^{\wedge})^{\mathcal{M}}} \arrow[dr,color=blue, "\hand{\approx}"] \& \& \hand{(B^{\wedge})_{\mathcal{M}}\overset{\mathrm{def}}{=}}\underline{\mathrm{Hom}}^{\hand{!}}(B^{\wedge\circ}\hand{,}\hand{\mathcal{M}}) \arrow[dl,color=blue, "\hand{\approx}"'] \\
(40) \& \underline{\mathrm{Hom}}(A,\hand{\mathcal{M}}) = \underline{\mathrm{Hom}}(B^{\circ},\hand{\mathcal{M}}) \& ,
\end{tikzcd}}

and depends upon the choice of a quasi-inverse of the second vertical
equivalence in (40) (which choice can be made, via the dual of (35),
via the choice of small inverse limits in \hand{$\mathcal{M}$}).

\uline{R\unsure{\over[typed]{g}{e}}mark.} It is felt that the way we got the equivalence (39) via (40),
the rôle of A and B in it is symmetric. To give a more precise statement,
consider the equivalence deduced from (39) by passing to the dual catego-%
ries of the two members - using the tautological isomorphisms (6) of
page 385, we get an equivalence \deleted{\hand{\unsure{$(\delta_{\illegible{1}}^{\mathcal M})^\circ$}}}
\[ (39')\quad \hand{(\delta^{\mathcal{M}}_{A})^{\circ}}:\ \underline{\mathrm{Hom}}^{\hand{!}}(A^{\wedge\circ},\hand{\mathcal{N}}) \hand{\xrightarrow{\approx}} \underline{\mathrm{Hom}}_{\hand{!}}(B^{\wedge},\hand{\mathcal{N}}) \quad , \quad \text{\over[typed]{w}{W}here } \hand{\mathcal{N}=\mathcal{M}^{\circ}}\ ; \]
th\unsure{\over[typed]{i}{e}}s equivalence is \uline{canonicall\unsure{\illegible{1}}y} \uline{quasi-inverse} \uline{to the equivalence} \unsure{\illegible{1}}\xout{\illegible{1}}
\hand{$\delta^{\mathcal{N}}_{B}$} \uline{in opposite direction}, associated to the pair (B,\hand{$\mathcal{N}=\mathcal{M}^{\circ}$}) instead of
(A,\hand{$\mathcal{M}$}) .

In the rest of this subsection D), we'll elaborate on some particu-%
lar cases of the eq\unsure{\over[typed]{m}{u}}ivalence (39) between cosheaves and sheaves.

\inserted{\uline{Case}}\over{1}{1}\textsuperscript{\uline{o}}). Assume \hand{$\mathcal{M}$} = (Sets), then by prop.1 (p.384) the right hand
side of (39) is canonically equivalent to $B^{\wedge}$ itself, hence we get an
equivalence
\[ (41)\qquad B^{\wedge} \hand{\xrightarrow{\approx}} \underline{\mathrm{Hom}}_{\hand{!}}(A^{\wedge},\ (\mathrm{Sets})) \quad . \]
If we want to keep track of the symmetry aspect described in the remark
above, we may consider \xout{(41) as} the functor (41) as being deduced from a
canonical"pairing" between the categories $A^{\wedge}$ and $B^{\wedge}$
\[ (42)\qquad \hand{\delta_{A}}:\ A^{\wedge}\ \mathrm{x}\ B^{\wedge} \hand{\longrightarrow} (\mathrm{Sets}) \hand{,} \]
which is an object in
\[ \underline{\mathrm{Hom}}_{\hand{!!}}(A^{\wedge},\ B^{\wedge};\ (\mathrm{Sets})) \]
i.e. which commutes to small direct limits in e ach variable. (For this
interpretation, compare the dual statement contained in f\unsure{\over[typed]{b}{o}}rmula (7) of
page 385 - and note that (41), being an equivalence, commutes to
small direct limits i.e. is in a category $\underline{\mathrm{Hom}}_{\hand{!}}(B^{\wedge},\ \underline{\mathrm{Hom}}_{\hand{!}}(\ \ ,\ \ ))$ .)

\margin{78}
